%======================================================================
\chapter{Introduction}\label{ch:intro}
%======================================================================

\section{What this volume is about}

The classification of homogeneous plane continua (\vone{Theorem~0.1}) splits into an indecomposable and a decomposable case. Volume~6
treated the indecomposable case. This volume treats the decomposable case:

\begin{theorem}[Theorem~\ref{thm:decomposable}]
Let $M$ be a decomposable homogeneous plane continuum. Then either $M$ is a simple closed curve, or there is a continuous, open, surjective map
$f\colon M\to\Sph$ whose fibres are pseudo-arcs and are terminal in $M$ (every continuum meeting a fibre lies in it or contains it).
\end{theorem}

The proof follows two papers of F.~B.~Jones.
\begin{enumerate}
\item \textbf{Jones's decomposition} (1955, Theorem~1). For $x\in M$ let $\Ls x$ be the set of points $z$ such that every continuum containing a
neighbourhood of $z$ contains $x$. In a homogeneous continuum the sets $\Ls x$ partition $M$ into terminal continua; the decomposition is continuous;
and the quotient $M/G$ is a homogeneous \emph{aposyndetic} continuum. When $M$ lies in the plane, the elements do not separate the plane, and $M/G$
is again a plane continuum.
\item \textbf{Jones's theorem of 1949.} A nondegenerate homogeneous aposyndetic plane continuum is a simple closed curve. So $M/G$ is a circle
(Jones 1955, Theorem~2).
\item \textbf{The elements.} Then $M$ has exactly two complementary domains, so a decomposable element would separate the plane; hence the elements are
indecomposable, and by Volume~6 the nondegenerate ones are pseudo-arcs.
\end{enumerate}
The circle-like structure of $M$, and the fact that there is only one circle of pseudo-arcs up to homeomorphism, are the subject of Volume~8.

The volume has nine chapters after this introduction.

\begin{itemize}
\item \textbf{Chapter~\ref{ch:apos}: Aposyndesis.} Jones's 1948 paper: openness, the sets $\Ls y$ are continua, semi-local connectedness, decomposability.
\item \textbf{Chapter~\ref{ch:spread}: Homogeneity and aposyndesis.} Spreading (from the Effros theorem), the Lemma of Jones 1951, and Bing's theorem on non-cutting points.
\item \textbf{Chapter~\ref{ch:jones55}: Jones's aposyndetic decomposition.} Theorem~1 of 1955, continuity of the decomposition, the elements, non-separation of the plane;
Jones's 1951 theorem on hereditarily unicoherent continua.
\item \textbf{Chapter~\ref{ch:order}: Order and Menger's theorem.} Order of a point; a continuum all of whose points have order~$2$ is a circle.
\item \textbf{Chapter~\ref{ch:whyburn}: Whyburn's separating point--order theorem.} Only countably many local separating points fail to have order~$2$, by chains of separating points.
\item \textbf{Chapter~\ref{ch:sectors}: Continua crossing an annulus.} The sector theorem and the cyclic order, by lifting to a strip; crossing continua from a failure of local connectedness.
\item \textbf{Chapter~\ref{ch:boundaries}: Boundaries of complementary domains.} Whyburn's 1939 Theorem~14 and its Corollary~1: for a semi-locally connected plane continuum the
boundaries of complementary domains are locally connected, and simple closed curves if there are no cut points.
\item \textbf{Chapter~\ref{ch:jones49}: Homogeneous aposyndetic plane continua are circles.} Jones 1949, Theorems~1 and~2.
\item \textbf{Chapter~\ref{ch:decomposable}: The decomposable case.} The plane of elements, Jones 1955 Theorem~2 and its remarks, the elements are pseudo-arcs, the main theorem.
\end{itemize}

The main results are:

\begin{theorem}[Jones 1955, Theorem~1; Theorems~\ref{thm:jones55-b}, \ref{thm:jones55-c}, \ref{thm:jones55-usc}, \ref{thm:jones55-a}]
In a homogeneous continuum $X$ the sets $\Ls x$ partition $X$ into terminal subcontinua; the decomposition is continuous; and if $X$ is decomposable, the quotient is a
nondegenerate homogeneous aposyndetic continuum.
\end{theorem}

\begin{theorem}[Menger; Theorem~\ref{thm:menger}]
A nondegenerate continuum all of whose points have order $2$ is a simple closed curve.
\end{theorem}

\begin{theorem}[Whyburn; Theorem~\ref{thm:sep-order}, Corollary~\ref{cor:whyburn}]
In a separable locally compact metric space all but countably many local separating points have order $2$.
\end{theorem}

\begin{theorem}[Whyburn 1939; Theorems~\ref{thm:whyburn14} and~\ref{thm:whyburn-cor1}]
The boundary of each complementary domain of a semi-locally connected plane continuum is locally connected; if the continuum has no cut point (and is nondegenerate), the
boundary is a simple closed curve.
\end{theorem}

\begin{theorem}[Jones 1949; Theorem~\ref{thm:jones49-1}]
A nondegenerate homogeneous aposyndetic plane continuum is a simple closed curve.
\end{theorem}

\begin{theorem}[Jones 1955, Theorem~2; Theorem~\ref{thm:jones55-2}, Corollary~\ref{cor:jones55-remarks}]
If $M$ is a decomposable homogeneous plane continuum, $M/G$ is a simple closed curve, each element is a homogeneous indecomposable continuum not separating the plane, and $M$
has exactly two complementary domains.
\end{theorem}

\section{Where these results are used}

Theorem~\ref{thm:decomposable} is the decomposable case of the classification, assembled in Volume~8 together with the indecomposable case (\vsix{Theorem~7.4}) and the
theorem of Bing and Jones that circles of pseudo-arcs are unique. The map $f$ of Theorem~\ref{thm:decomposable} is the starting point of Volume~8.

\section{Where the formal proof differs from the literature}\label{sec:intro-diff}

\begin{enumerate}
\item \textbf{Jones 1955.} Parts (b) and (c) of Theorem~1 are proved from a consequence of the Effros theorem (spreading, Theorem~\ref{thm:spreading}) and the Lemma of Jones
1951, proved with a lower semicontinuous function in place of a transfinite sequence; Jones's twelve claims and his case argument are not needed (they are proved formally
nonetheless). Decomposability is not needed for (b) and (c). Aposyndesis of the quotient, which Jones states without argument, is proved from the continuity of the
decomposition. In the non-separation argument the disjointness of the domains $W_g$ is supplied (Theorem~\ref{thm:elt-not-separates}). The indecomposability of the elements is
derived from the two-domain theorem rather than from Jones's 1951 paper.
\item \textbf{Bing 1948.} Bing's Theorem~5 on non-cutting points is proved by a sequence of tight continua (Section~\ref{sec:bing5}).
\item \textbf{Menger and Whyburn.} Menger's theorem is proved in the special case needed, by a triod argument. Whyburn's separating point--order theorem is proved with chains
of separating points and a cut-wire argument, without his theory of cyclic elements and potential order.
\item \textbf{Moore 1922.} The two facts about crossing continua that Moore takes for granted (the sector theorem and the cyclic order) are proved by lifting to a strip, with
continuous logarithms only.
\item \textbf{Whyburn 1939.} The reduction to bounded domains by inversion has a gap and is avoided; Whyburn's arcs, his use of a theorem of Moore and his curve through a
point of order~$2$ are replaced by Janiszewski's theorem, the sector theorem and accumulating continua. Corollary~1, stated without proof, is proved with a $\theta$-curve lemma.
\item \textbf{Jones 1949.} Theorems~1 and~2 are false as printed for a one-point continuum. ``$M-(A+B)$ is not connected'' is proved with the $\theta$-curve lemma
(Lemma~\ref{lem:theta-sep}), and the two cases end with explicit local separating points.
\end{enumerate}

\section{Prerequisites and conventions}\label{sec:intro-conv}

The reader should know the material of Volume~3 (boundary bumping, irreducible continua, non-cut points, upper semicontinuous decompositions and their continuity at a
residual set of elements, the Effros theorem), the plane topology of Volume~4 (continuous logarithms, Janiszewski's theorem, Brouwer's theorem, invariance of plane separation)
and Volume~5 (arcs and simple closed curves, the Jordan curve theorem, Moore's theorem on decompositions of the plane). Chapter~\ref{ch:decomposable} uses the main theorem of
Volume~6.

\begin{itemize}
\item \textbf{Continua.} A continuum is compact, connected and metric; it may be a point unless ``nondegenerate'' is said.
\item \textbf{The plane.} As in Volumes~4 to~6, $\R^2$ carries the sup metric, so that balls are squares.
\item \textbf{Simple closed curves.} A simple closed curve is the union of two arcs with the same end points and no other common point; by \vfive{Proposition~1.13} it is
homeomorphic to the circle $\Sph$.
\end{itemize}

\paragraph{Formal counterparts.} As in the earlier volumes, many results are followed by a line ``Formal counterpart:'' naming the corresponding declaration of the
Lean development (all names begin with \lean{HPC.}). These lines may be skipped. Appendix~\ref{app:formal} collects them.

\paragraph{Exercises.} Each chapter ends with exercises. None of them is used later.
